\subsection{Criterios de diseño}

El programa se ha diseñado siguiendo una arquitectura modular, de forma que cada archivo tenga una responsabilidad concreta y no se mezclen decisiones de alto nivel con detalles eléctricos. Esta separación resulta especialmente importante en un prototipo electromecánico, ya que permite modificar un pin, ajustar un tiempo de rebote o cambiar un umbral de la botonera sin alterar la lógica general de reparto.

Los criterios principales han sido los siguientes:

\begin{itemize}
    \item \textbf{Determinismo:} el sistema debe responder siempre de la misma forma ante la misma secuencia de eventos.
    \item \textbf{No bloqueo:} el código evita esperas con \texttt{delay()} y trabaja mediante actualizaciones periódicas basadas en \texttt{millis()}.
    \item \textbf{Separación por capas:} la FSM decide qué debe ocurrir, pero no manipula pines directamente.
    \item \textbf{Seguridad de operación:} los movimientos de giro y empuje incorporan \textit{timeouts} para evitar que un motor quede funcionando indefinidamente ante un fallo de sensor.
    \item \textbf{Calibración centralizada:} pines, umbrales, tiempos y velocidades PWM se agrupan en \texttt{src/config.h}.
\end{itemize}

\subsection{Estructura del repositorio}

El punto de entrada del programa es mínimo. El archivo \texttt{card\_dealer.ino} únicamente inicializa la máquina de estados en \texttt{setup()} y la actualiza de forma continua en \texttt{loop()}. El comportamiento real se delega en los módulos de \texttt{src/}.

\begin{table}[H]
\centering
\small
\begin{tabularx}{\textwidth}{>{\bfseries}p{0.30\textwidth}X}
\hline
\textbf{Archivo} & \textbf{Responsabilidad} \\
\hline
\texttt{card\_dealer.ino} & Entrada principal de Arduino. Llama a \texttt{fsm\_init()} y \texttt{fsm\_update()}. \\
\texttt{src/config.h} & Define pines, límites de jugadores, tiempos de seguridad, PWM, polaridad de sensores y umbrales de la botonera. \\
\texttt{src/config\_checks.h} & Verifica en compilación que no existan duplicados de pines ni incoherencias básicas de configuración. \\
\texttt{src/fsm.*} & Implementa la lógica del sistema mediante una máquina de estados finitos. Decide cuándo girar, empujar cartas, volver a HOME o entrar en error. \\
\texttt{src/hardware.*} & Centraliza el acceso a motores DC, drivers L298N, sensor Hall HOME y sensores digitales de carta. \\
\texttt{src/ui.*} & Gestiona la pantalla LCD 16x2 y la botonera analógica conectada a A0. Traduce pulsaciones físicas en eventos lógicos. \\
\texttt{src/diagnostics.*} & Permite probar sensores y motores de forma independiente durante el montaje y calibración. \\
\texttt{src/debug.*} & Encapsula los mensajes por \texttt{Serial} mediante macros de depuración. \\
\hline
\end{tabularx}
\caption{Organización del software embebido del ELCO-Dealer}
\label{tab:software_files}
\end{table}

\subsection{Mapa de pines y elementos controlados}

El mapa de pines se encuentra en \texttt{Config} y actúa como única fuente de verdad. Esta decisión evita que aparezcan números de pin dispersos por el código y facilita la revisión del cableado.

\begin{table}[H]
\centering
\small
\begin{tabularx}{\textwidth}{>{\bfseries}p{0.38\textwidth}X}
\hline
\textbf{Elemento} & \textbf{Pines configurados} \\
\hline
LCD 16x2 & D4--D7 en 4, 5, 6 y 7; RS en 8; EN en 9; retroiluminación en 10. \\
Botonera LCD Keypad & Entrada analógica A0, interpretada mediante umbrales de tensión. \\
L298N 1 & PWM en 11 y 12; dirección en 22, 23, 24 y 25. \\
L298N 2 & PWM en 2 y 3; dirección en 26, 27, 28 y 29. \\
Sensores de carta & \texttt{CARD\_SENSOR\_1}, \texttt{CARD\_SENSOR\_2} y \texttt{CARD\_SENSOR\_3} en los pines 30, 31 y 32. \\
Sensor HOME & Sensor Hall en el pin 33. \\
\hline
\end{tabularx}
\caption{Resumen del mapa de pines definido en \texttt{src/config.h}}
\label{tab:pin_map}
\end{table}

La versión actual del software utiliza tres sensores de carta y tres canales de reparto. Las parejas \texttt{DEAL\_1}/\texttt{CARD\_1}, \texttt{DEAL\_2}/\texttt{CARD\_2} y \texttt{DEAL\_3}/\texttt{CARD\_3} se definen en \texttt{src/config.h}.

\subsection{Máquina de estados finitos}

La FSM es el núcleo del sistema. Su función es coordinar la secuencia de operación del prototipo sin acceder directamente a pines ni duplicar lógica de hardware. Cada transición se produce por un evento de usuario, por la detección de un flanco de sensor o por un \textit{timeout} de seguridad.

\begin{table}[H]
\centering
\small
\begin{tabularx}{\textwidth}{>{\bfseries}p{0.30\textwidth}X}
\hline
\textbf{Estado} & \textbf{Función dentro del sistema} \\
\hline
\texttt{OFF} & Estado seguro inicial. Todos los motores están detenidos y la pantalla indica que se mantenga pulsado SELECT. \\
\texttt{INIT} & Inicialización lógica tras encendido. Reinicia contadores y prepara el menú de configuración. \\
\texttt{SET\_PLAYERS} & Selección del número de jugadores, con límite máximo de 6. \\
\texttt{SET\_CARDS} & Selección del número de cartas por jugador, con límite definido por \texttt{CMAX}. \\
\texttt{AUTO\_INIT} & Preparación del reparto automático. Reinicia el progreso de la ronda. \\
\texttt{ROTATE\_AUTO} & Gira la plataforma hasta el slot objetivo del jugador actual. \\
\texttt{PUSH\_AUTO} & Acciona el motor de reparto del canal activo hasta detectar una carta mediante su sensor asociado. \\
\texttt{RETURN\_HOME\_AUTO} & Devuelve el sistema a la posición HOME tras completar la ronda automática. \\
\texttt{IDLE\_HOME} & Estado de reposo operativo. Permite lanzar otra ronda o repartir una carta manual. \\
\texttt{ROTATE\_ONE} & Gira hacia el siguiente jugador en modo manual. \\
\texttt{PUSH\_ONE} & Reparte una única carta y espera su detección. \\
\texttt{RETURN\_HOME\_MANUAL} & Vuelve a HOME tras el reparto manual. \\
\texttt{ERROR} & Estado seguro ante fallo de giro, carta o retorno a HOME. Los motores se detienen. \\
\hline
\end{tabularx}
\caption{Estados principales implementados en \texttt{src/fsm.cpp}}
\label{tab:fsm_states}
\end{table}

La distribución de jugadores se calcula sobre 6 posiciones lógicas, correspondientes a los imanes o slots de referencia. Para cada jugador, la función de cálculo asigna un slot objetivo proporcional al número total de jugadores configurado:

\[
\text{slot objetivo} = \left\lfloor \frac{\text{índice de jugador} \cdot 6}{\text{número de jugadores}} \right\rfloor
\]

Con esta estrategia se mantiene una separación uniforme entre jugadores y se evita invertir el sentido de giro. La plataforma gira siempre en el mismo sentido, y el primer imán detectado durante un retorno a HOME redefine el origen lógico como slot 0.

\subsection{Gestión de eventos de usuario}

La interfaz física se basa en la botonera de una LCD Keypad Shield cableada manualmente. El módulo \texttt{ui.cpp} lee el valor analógico de A0 y lo transforma en botones físicos. En la versión actual se utilizan dos acciones lógicas:

\begin{itemize}
    \item \textbf{PWR:} asociado al botón SELECT. Una pulsación corta confirma o avanza; una pulsación larga enciende o apaga lógicamente el sistema.
    \item \textbf{ADD:} asociado al botón RIGHT. Incrementa jugadores, cartas o lanza una carta manual según el estado.
\end{itemize}

La detección de pulsaciones utiliza un tiempo de antirrebote de \texttt{UI\_DEBOUNCE\_MS} y diferencia pulsación corta y larga mediante \texttt{UI\_LONG\_PRESS\_MS}. Internamente, los eventos se priorizan para que \texttt{PWR\_LONG} tenga precedencia sobre eventos de menor importancia.

\subsection{Control de sensores}

Los sensores digitales se gestionan exclusivamente desde \texttt{hardware.cpp}. Cada sensor mantiene su propio estado estable, último valor leído, instante del último cambio, instante del último evento y banderas de flanco ascendente o descendente.

El sensor Hall HOME utiliza:

\begin{itemize}
    \item lectura digital con nivel activo configurable en \texttt{HOME\_SENSOR\_ACTIVE\_LEVEL};
    \item antirrebote mediante \texttt{HALL\_DEBOUNCE\_MS};
    \item tiempo muerto mediante \texttt{HALL\_DEAD\_TIME\_MS};
    \item rearme lógico en la FSM mediante \texttt{HALL\_REARM\_MS}.
\end{itemize}

El sensor de carta sigue la misma filosofía, pero con tiempos específicos definidos como \texttt{CARD\_SENSOR\_DEBOUNCE\_MS}, \texttt{CARD\_SENSOR\_DEAD\_TIME\_MS} y \texttt{CARD\_SENSOR\_REARM\_MS}. De esta forma se evita contar varias veces una misma carta cuando el sensor permanece activo durante el paso físico del naipe.

\subsection{Programación de los motores}

Todos los motores DC se controlan desde la capa \texttt{Hardware}. La FSM no escribe pines directamente: solicita acciones lógicas como \texttt{motor\_forward()}, \texttt{motor\_stop()} o \texttt{motors\_stop\_all()}.

Cada canal de motor se representa mediante:

\begin{itemize}
    \item un pin PWM conectado al enable del L298N;
    \item dos pines de dirección;
    \item una velocidad PWM actual;
    \item una dirección lógica.
\end{itemize}

El motor de rotación se controla como \texttt{MotorId::ROTATION} y los motores de reparto como \texttt{DEAL\_1}, \texttt{DEAL\_2} y \texttt{DEAL\_3}. La FSM selecciona el canal de reparto según el jugador actual y confirma la carta con el sensor asociado a ese canal.

La velocidad inicial de giro se define en \texttt{ROTATION\_PWM} y la velocidad de reparto en \texttt{DEAL\_PWM}. Estos valores deben calibrarse sobre el prototipo real, ya que dependen de la fricción de los rodillos, el peso de la baraja, la alimentación de los motores y la alineación del mecanismo.

\subsection{Timeouts y recuperación de errores}

Para que el sistema sea tolerante a fallos físicos, las acciones con movimiento tienen un tiempo máximo de ejecución:

\begin{table}[H]
\centering
\small
\begin{tabularx}{\textwidth}{>{\bfseries}p{0.35\textwidth}X}
\hline
\textbf{Parámetro} & \textbf{Uso} \\
\hline
\texttt{ROTATE\_TIMEOUT\_MS} & Tiempo máximo permitido para alcanzar un slot durante el giro. \\
\texttt{PUSH\_TIMEOUT\_MS} & Tiempo máximo para detectar una carta tras activar el motor de reparto. \\
\texttt{RETURN\_HOME\_TIMEOUT\_MS} & Tiempo máximo para volver a detectar HOME durante el retorno. \\
\hline
\end{tabularx}
\caption{Timeouts principales de seguridad}
\label{tab:timeouts}
\end{table}

Si se supera alguno de estos límites, la FSM entra en \texttt{ERROR}, detiene todos los motores y muestra el tipo de fallo en la pantalla. Desde este estado se puede iniciar una recuperación mediante una pulsación corta de PWR, que ordena un retorno a HOME.

\subsection{Modo de diagnóstico}

El modo de diagnóstico permite validar el prototipo antes de ejecutar la FSM completa. Se activa en compilación mediante la macro:

\begin{verbatim}
CARD_DEALER_DIAGNOSTIC_MODE=true
\end{verbatim}

En este modo, el usuario navega por páginas de prueba usando la pantalla LCD:

\begin{itemize}
    \item \textbf{PWR corto:} avanzar a la siguiente página de diagnóstico.
    \item \textbf{ADD corto:} activar o detener el motor mostrado en pantalla.
    \item \textbf{PWR largo:} parada segura y vuelta a la página HOME.
\end{itemize}

El recorrido recomendado es comprobar primero el sensor HOME, después los tres sensores de carta y finalmente cada motor. No debe pasarse a modo normal hasta que los flancos de sensor sean repetibles y ningún motor quede activado al cambiar de pantalla.

\subsection{Compilación y carga}

El proyecto está preparado para compilarse con Arduino CLI usando la placa \textbf{Arduino Mega}. En modo normal:

\begin{verbatim}
arduino-cli compile --fqbn arduino:avr:mega .
\end{verbatim}

Para compilar el modo de diagnóstico sin modificar el archivo de configuración:

\begin{verbatim}
arduino-cli compile --fqbn arduino:avr:mega `
  --build-property compiler.cpp.extra_flags=-DCARD_DEALER_DIAGNOSTIC_MODE=true .
\end{verbatim}

La carga en placa debe realizarse indicando el puerto serie real detectado en el equipo:

\begin{verbatim}
arduino-cli upload -p COMx --fqbn arduino:avr:mega .
\end{verbatim}

Antes de cargar una versión de funcionamiento normal es recomendable revisar que \texttt{CARD\_DEALER\_DIAGNOSTIC\_MODE} esté desactivado o que la compilación no esté sobrescribiendo dicha macro.

\subsection{Parámetros de calibración}

La calibración final debe hacerse con el prototipo físico montado. Los parámetros más sensibles se encuentran en \texttt{src/config.h}:

\begin{table}[H]
\centering
\small
\begin{tabularx}{\textwidth}{>{\bfseries}p{0.32\textwidth}X}
\hline
\textbf{Grupo} & \textbf{Parámetros a revisar} \\
\hline
Botonera & Umbrales \texttt{KEYPAD\_RIGHT\_MAX}, \texttt{KEYPAD\_UP\_MAX}, \texttt{KEYPAD\_DOWN\_MAX}, \texttt{KEYPAD\_LEFT\_MAX} y \texttt{KEYPAD\_SELECT\_MAX}. \\
Sensores & Polaridad \texttt{CARD\_SENSOR\_ACTIVE\_LEVEL} y \texttt{HOME\_SENSOR\_ACTIVE\_LEVEL}. \\
Filtrado & Tiempos de antirrebote, tiempo muerto y rearme para Hall y sensores de carta. \\
Motores & \texttt{ROTATION\_PWM} y \texttt{DEAL\_PWM}. \\
Seguridad & \texttt{ROTATE\_TIMEOUT\_MS}, \texttt{PUSH\_TIMEOUT\_MS} y \texttt{RETURN\_HOME\_TIMEOUT\_MS}. \\
\hline
\end{tabularx}
\caption{Parámetros software que requieren calibración física}
\label{tab:calibration}
\end{table}

Un ajuste correcto debe conseguir que cada sensor de carta genere un único evento por carta, que HOME se detecte una sola vez por imán y que los motores tengan margen suficiente para mover la mecánica sin golpes bruscos.

\subsection{Flujo de uso desde software}

El flujo normal del programa se puede resumir del siguiente modo:

\begin{enumerate}
    \item El sistema arranca en \texttt{OFF} con todos los motores detenidos.
    \item Una pulsación larga de SELECT pasa a \texttt{INIT}.
    \item En \texttt{SET\_PLAYERS}, RIGHT incrementa el número de jugadores y SELECT confirma.
    \item En \texttt{SET\_CARDS}, RIGHT incrementa las cartas por jugador y SELECT inicia el reparto automático.
    \item En reparto automático, la FSM alterna entre giro y empuje hasta completar todas las cartas.
    \item Al terminar, el sistema vuelve a HOME y queda en \texttt{IDLE\_HOME}.
    \item Desde \texttt{IDLE\_HOME}, SELECT repite un reparto automático y RIGHT reparte una carta individual.
\end{enumerate}

Este flujo coincide con la experiencia prevista para el usuario final, pero mantiene internamente los controles de seguridad y recuperación necesarios para un prototipo físico.
